\documentclass[UTF8,fontset=macnew,11pt]{ctexart}

\usepackage[a4paper,margin=2.25cm]{geometry}
\usepackage{amsmath,amssymb}
\usepackage{booktabs}
\usepackage{tabularx}
\usepackage{array}
\usepackage{float}
\usepackage{xcolor}
\usepackage{hyperref}

\hypersetup{
  colorlinks=true,
  linkcolor=blue!50!black,
  citecolor=blue!50!black,
  urlcolor=blue!50!black
}

\setlength{\parskip}{0.34em}
\setlength{\parindent}{2em}
\setlength{\emergencystretch}{2em}

\newcommand{\method}{Resource-NMCTS}
\newcommand{\anf}{\ensuremath{\mathrm{ANF}}}
\newcommand{\fprm}{\ensuremath{\mathrm{FPRM}}}
\newcommand{\mcts}{\ensuremath{\mathrm{MCTS}}}
\newcommand{\xor}{\ensuremath{\oplus}}
\newcommand{\mct}{\ensuremath{\mathrm{MCT}}}
\newcommand{\rz}{\ensuremath{R_z}}
\newcommand{\score}{\ensuremath{\mathrm{score}}}
\newcommand{\ttable}{truth-table}

\title{基于强化学习与蒙特卡洛树搜索的量子布尔函数综合方法：资源约束搜索视角}
\author{中文论文稿 v33：n=24 完整验证桥接版}
\date{2026年7月8日}

\begin{document}
\pagestyle{plain}
\maketitle

\begin{abstract}
量子布尔函数 oracle 综合把经典谓词嵌入量子算法，是 Grover 搜索、量子密码分析和约束求解类量子程序中的基础编译环节。现有方法通常从固定布尔表示、可逆模板、LUT 映射、XAG 网络或 CNOT 导向空间结构出发；这些路线能够在局部指标上取得效果，但很难同时处理 T-count、CNOT、深度、门数和辅助比特之间的组合权衡。本文将 Boolean oracle synthesis 重新表述为代数正规形（algebraic normal form, \anf{}）和固定极性 Reed--Muller（fixed-polarity Reed--Muller, \fprm{}）项集合上的资源约束搜索问题，并提出 \method{}：一种结合神经动作先验、蒙特卡洛树搜索、布尔环线性因子筛选、结构深度策略、depth-frontier policy 与保守 guard 的逻辑层综合框架。本文方法不依赖 SSHR 的 parallelotope 候选空间；SSHR 仅作为 small-scale CNOT-oriented baseline。

在 $n\leq6$ 的 177 个传统函数上，Pareto-\method{} 相对 direct \anf{} 平均 T-count 降低 72.25\%，平均 score 降低 67.80\%；相对 ESOP cube beam 获得 174/0/3 的 score 胜/负/平，平均 score 降低 36.09\%；相对 weighted ESOP-MILP 获得 167/3/7，平均 score 降低 29.84\%。在外部工具链 probe 中，ROS-style LUT proxy 覆盖 309 个函数和 927 个 ABC 映射行，best-$K$ proxy 相对固定 $K=4$ 平均 score 降低 18.12\%，但 Pareto-\method{} 相对该更强 proxy 仍为 309/0/0，平均 score 降低 84.27\%。官方 mockturtle KLUT-to-XAG probe 在传统 177 个函数上给出 166/11/0、平均 score 降低 31.50\%，在 $n=14$ 的 64 个高维函数上给出 64/0/0、平均 score 降低 91.49\%。高维正确性通过两层桥接补强：$n=20$--$40$ 项集级 6600/6600 个方法行通过 \anf{} plan 展开验证和 emitted circuit 符号验证；$n=21,22,23$ truth-table bridge 中 300/300 个方法行通过完整 oracle 验证。本文进一步用位掩码 truth-table evaluator 将完整 bridge 推进到 $n=24$：在 6 个生成式 \anf{} 函数上，60/60 个方法行同时通过完整 oracle、plan \anf{} 和 emitted-circuit 符号验证，平均 truth-table 构造时间为 0.18 s/函数。

本文同时报告 gate-set 边界。RevKit Python \texttt{oracle\_synth} 在 177 个函数上全部返回，但 171/177 行含非 Clifford \rz{}，总数为 9242，因此其 lower-bound phase netlist 不能直接当作精确 Clifford+T T-count 对比。为把该边界推进到项目内部，本文实现 phase-parity \anf{} emitter：将每个 \anf{} 单项式精确展开为 parity-phase gadgets，177/177 个函数均验证为 phase oracle up to global phase。该 emitter 相对 RevKit lower-bound score 为 40/137/0、平均高 69.25\%，但在每个非 Clifford \rz{} 加 1 个 score 单位后转为 177/0/0、平均降低 48.16\%，在 $T/R_z=30$ 近似旋转综合 proxy 下为 177/0/0、平均降低 64.98\%。进一步地，本文新增 fixed-polarity phase-parity search：对每个函数枚举 \fprm{} 极性 $p$，综合 $g(z)=f(z\xor p)$ 的 phase polynomial，再把 shifted parity 精确翻译回原变量。三个 rank metric 共 531 行全部验证通过；其中 $T/R_z=30$ 目标在 59/177 个函数选择非零极性，相对原始 phase-parity \anf{} 为 59/0/118、平均 score 降低 0.47\%，相对 RevKit $T/R_z=30$ proxy 为 177/0/0、平均降低 65.16\%。这说明 \fprm{} 极性已经把 phase/Rz 方向变成实际搜索问题，但收益仍小，不能把它写成最终 phase 方法。阶段性结论是：\method{} 已经形成一条独立于 SSHR、以低 T-count 与低加权逻辑资源为目标的 Boolean oracle synthesis 主线；当前缺口是更强的学习式 phase/Rz-aware search、官方 ROS 或 legacy RevKit/CirKit 完整复现，以及真实后端 mapping 评估。
\end{abstract}

\noindent\textbf{关键词：}量子 oracle 综合；布尔函数综合；强化学习；蒙特卡洛树搜索；ANF；FPRM；资源约束编译

\section{本文定位：把 oracle 综合写成搜索问题}

本文的核心问题是：给定布尔函数 $f:\{0,1\}^n\rightarrow\{0,1\}$，如何自动生成实现
\begin{equation}
  |x\rangle |y\rangle \mapsto |x\rangle |y\xor f(x)\rangle
  \label{eq:oracle}
\end{equation}
的逻辑层可逆线路，并在 T-count、CNOT、深度、门数和辅助比特之间取得更好的资源权衡。这个问题本质上不是单一规则匹配问题，而是组合搜索问题：同一个函数可以有许多等价的布尔表示、因子分解、极性选择和计算/反计算顺序，不同选择会导致完全不同的量子逻辑资源。

本文不从 SSHR 入手。SSHR 的贡献在于利用 parallelotope 空间结构做 small-scale CNOT-oriented synthesis\cite{zheng2025sshr}，适合作为 CNOT 导向 baseline；但本文要解决的是更一般的资源约束搜索问题。换言之，SSHR 在本文中回答的是“CNOT-oriented small-scale baseline 怎么样”，而不是决定本文方法空间。本文的状态、动作、验证和资源模型都定义在 \anf{}/\fprm{} 项集合及其代数分解上。

本文主张可以压缩为一句话：

\begin{quote}
在逻辑层 bit-flip Boolean oracle synthesis 中，用 \anf{}/\fprm{} 项集合表示状态、用可反算的代数分解定义动作，并用神经先验和 \mcts{} 控制搜索，可以在传统函数、高维项集合和若干外部逻辑网络 probe 上降低 T-count 与加权 score；该主张不覆盖 phase/Rz oracle 最优性、硬件 routing、magic-state factory、噪声模型或后端映射最优性。
\end{quote}

这个边界是本文能否投稿的关键。若把文章写成“SSHR 的改进”，贡献会被限制在 small-scale CNOT-oriented synthesis；而从搜索问题出发，本文可以自然比较 SSHR、ESOP、CNOT、T-count、辅助比特、线路长度、深度和 phase/Rz gate-set 边界。AI 的角色也更清楚：它不是直接生成无约束门序列，而是在可验证动作空间中分配搜索预算。

\begin{table}[H]
\centering
\small
\caption{本文固定术语与使用范围。}
\label{tab:terms}
\begin{tabularx}{\linewidth}{>{\raggedright\arraybackslash}p{0.27\linewidth}>{\raggedright\arraybackslash}X}
\toprule
术语 & 本文含义 \\
\midrule
\method{} & 在 \anf{}/\fprm{} 项集合上进行资源加权搜索的主框架，包括神经动作先验、\mcts{}、结构策略、guard 和 portfolio 选择。 \\
Boolean-ring screen & 在 square-free Boolean ring 中搜索可复用线性因子的结构筛选分支，允许线性因子与 quotient 共享变量。 \\
Depth-frontier policy & 在 depth-2、depth-3 和 depth-4 screen 之间选择的高预算质量-时间折中策略。 \\
Guard & 以确定性 baseline 为兜底的保守门控；只有学习候选不劣于 baseline 时才接受。 \\
Truth-table bridge & 在 $n=21,22,23,24$ 构造完整 \ttable{}，把大规模项集符号验证和函数级完整验证连接起来。 \\
SSHR & CNOT-oriented small-scale baseline；本文方法不使用 SSHR parallelotope 候选空间。 \\
\bottomrule
\end{tabularx}
\end{table}

\section{相关工作与差异}

\subsection{布尔表示与低 T-count oracle 综合}

低 T-count oracle 综合长期依赖合适的布尔中间表示。Xor-And-Inverter Graphs（XAG）把 XOR 和 AND 结构分离，使 AND 节点数与非 Clifford 成本形成清晰联系。Meuli 等人讨论了 multiplicative complexity 在低 T-count oracle 编译中的作用\cite{meuli2019multiplicative}，并将 XAG 用于量子编译\cite{meuli2022xag}。ROS 将 oracle synthesis 写成 resource-constrained LUT mapping\cite{meuli2020ros}，进一步表明 LUT 分解、局部函数实现和辅助比特管理会共同影响资源。BDD-based reversible synthesis 展示了面向大函数的符号结构路线\cite{wille2009bdd}。后端感知 oracle synthesis 则开始把逻辑综合和 fault-tolerant back-end 指标连接起来\cite{yu2025backend}。

本文与这些工作的共同点是承认布尔中间表示决定量子逻辑资源。差异在于，本文不固定使用 XAG、LUT 或 BDD，而是以 \anf{}/\fprm{} 项集合为搜索状态。这样做的优势是任意候选 plan 都能展开回 GF(2) 多项式进行等价检查，且动作可以包含公共因子、极性重写、线性奇偶因子和布尔环共享变量结构。代价是需要重新设计搜索策略和外部对照解释，尤其要避免把逻辑层 proxy 写成硬件级结论。

\subsection{学习式搜索与线路综合}

学习式搜索已经进入经典电路设计、量子线路综合和图重写优化。ShortCircuit 使用 AlphaZero 风格搜索进行经典布尔电路设计\cite{tsaras2024shortcircuit}；Gumbel AlphaZero 被用于 Clifford+T unitary synthesis\cite{rietsch2024unitary}；强化学习和图神经网络也被用于 ZX-calculus rewrite-rule 选择\cite{riu2025rlzx}；MonteQ 使用 \mcts{} 处理量子线路综合中的动作排序问题\cite{machiya2026monteq}。这些研究说明，离散综合中的动作空间往往过大，学习策略可以帮助搜索更快进入高质量区域。

本文的状态和动作与上述工作不同。我们不让模型直接输出门序列，而是让模型在可验证的代数动作之间排序或选择。最终线路由 compute/action/uncompute 结构生成，并经过 \anf{} 展开、GF(2) 多项式模拟或完整 \ttable{} 验证。这样的设计降低了 AI 产生错误线路的风险，也使实验可以清楚区分“学习策略带来的搜索改进”和“确定性代数结构本身的收益”。

\section{问题定义与资源模型}

本文输入可以是完整 \ttable{}，也可以是 \anf{} 项集合。对于
\begin{equation}
  f(x)=\bigoplus_{m\in\mathcal{M}}\prod_{i\in m}x_i ,
  \label{eq:anf}
\end{equation}
$\mathcal{M}$ 表示 square-free monomial 集合。输出线路使用 X、CNOT 和多控 Toffoli（\mct{}）抽象门，并在逻辑层估计 T-count、CNOT、depth、gate count 和 peak ancilla。搜索排序使用统一分数
\begin{equation}
  \score = T + 0.04\,\mathrm{CNOT} + 0.015\,\mathrm{depth}
         + 0.01\,\mathrm{gates} + 2\,\mathrm{ancilla}.
  \label{eq:score}
\end{equation}
该分数只用于逻辑层候选排序，不代表真实硬件运行时间、物理错误率或 magic-state factory 总成本。因此，实验同时报告 T、CNOT、深度和辅助比特，避免单一 score 掩盖 trade-off。

正确性验证分为三层。第一层在 $n\leq6$ 传统函数、$n=18,20$ 函数级切片和 $n=21,22,23,24$ bridge 中构造完整 \ttable{}，逐点验证式~\eqref{eq:oracle}。第二层在大规模项集实验中，将 factor plan 展开回 \anf{} 项集合，检查目标多项式是否一致。第三层对 emitted X/CNOT/\mct{} circuit 做 GF(2) 多项式符号模拟，检查输出线多项式、输入线复原和辅助线清零。第三层不是 $2^n$ 枚举，但验证的是最终 emitted circuit，而不只是搜索 plan。

为使 $n=24$ 完整枚举可运行，本文在验证 harness 中加入位掩码 truth-table evaluator。对变量 $x_i$ 预先构造赋值位集 $V_i$，单项式 $m=\prod_{i\in S}x_i$ 的真值表为
\begin{equation}
  T_m=\bigwedge_{i\in S}V_i ,
  \label{eq:truth-mask}
\end{equation}
常数项对应全 1 掩码，完整 \anf{} 真值表则由所有 $T_m$ 异或得到。这个实现只改变完整验证的构造方式，不改变搜索算法或资源统计；小规模随机 \anf{} 已与原 zeta 实现逐项对比通过。

\section{方法}

\subsection{项集合搜索状态}

\method{} 的状态是当前尚未综合的项集合 $\mathcal{M}$。一次动作选择一个可计算、可反算的代数结构，把 $\mathcal{M}$ 分解为 quotient 子问题和 rest 子问题，再递归生成线路。候选动作包括公共单项式因子、\fprm{} 极性重写后的公共项、线性奇偶因子、布尔环线性因子和由神经模型扩展的 root actions。动作被接受后，框架会估计即时资源收益、子问题大小、辅助比特峰值和后续可分解性。

神经动作先验用于对候选排序，\mcts{} 在有限预算内探索候选组合，rollout 使用确定性资源估计和已有 baseline 候选。最终 portfolio 在多个候选线路中按式~\eqref{eq:score} 选择最优者。这个流程的关键不是“模型足够聪明”，而是每个候选动作都可以回到 GF(2) 语义验证；模型只影响搜索顺序和预算分配，不改变正确性定义。

\subsection{布尔环线性因子}

传统线性因子筛选通常要求线性因子变量和 quotient 变量不相交。这个限制便于实现，但会排除 square-free Boolean ring 中合法的重叠结构。本文的 Boolean-ring screen 允许二者共享变量，并在展开时使用 $x_i^2=x_i$ 与 GF(2) 偶数项消去。例如
\begin{equation}
  (x_i\xor x_j)x_i m = x_i m \xor x_i x_j m .
  \label{eq:boolean-ring}
\end{equation}
线路层面仍可用 CNOT 计算线性奇偶量，再以该辅助量控制 quotient 子线路。这样，screen 不只是把 beam 宽度调大，而是把原先被实现约束排除的代数候选重新放回可综合空间。

\subsection{结构策略、frontier policy 与 guard}

本文把 AI 模块拆成三类。第一类是动作排序，模型根据候选覆盖率、项数变化、变量覆盖密度、次数分布和即时资源估计对候选动作打分。第二类是结构选择，depth policy 判断当前项集合适合 single、depth-1 还是 depth-2 screen。第三类是 frontier 选择，depth-frontier policy 在 depth-2、depth-3 和 depth-4 screen 之间选择。质量型 frontier policy 用 score-only oracle frontier 训练，目标是在接近 all-depth frontier 的质量时减少全深度枚举开销；cost-aware frontier policy 则在标签中加入相对运行时间惩罚，用小幅 score 代价换取更低 plan time 和更低辅助线 lifetime area。

Guard 用来限制学习策略风险。若学习候选相对确定性 baseline 出现 score 退化，则返回 baseline。Depth-2 skip guard 避免把应该运行 depth-2 的项集合错误跳过。Screen-gate 用于判断某些边界函数是否可以跳过昂贵的 Resource tail。这样的设计让 AI 贡献可以表现为质量提升、运行时间节省或质量-时间折中，而不是要求学习策略在所有切片上无条件支配确定性 baseline。

\begin{table}[H]
\centering
\small
\caption{\method{} 的模块与证据钩子。}
\label{tab:modules}
\begin{tabularx}{\linewidth}{>{\raggedright\arraybackslash}p{0.24\linewidth}>{\raggedright\arraybackslash}X>{\raggedright\arraybackslash}p{0.24\linewidth}}
\toprule
模块 & 机制 & 证据钩子 \\
\midrule
\anf{}/\fprm{} 搜索状态 & 把 oracle 综合写成项集合分解问题，所有候选可回到 GF(2) 多项式验证。 & 小规模 \ttable{} 验证与项集级 \anf{} 验证。 \\
Boolean-ring screen & 搜索允许共享变量的线性因子，递归处理 quotient/rest 子问题。 & $n=18,20$ 函数级边界和 $n=20$--$40$ scale。 \\
神经动作先验与 \mcts{} & 对候选动作排序并在有限预算内组合分解动作。 & 传统函数上的搜索消融和 portfolio 改善。 \\
Depth-frontier policy & 学习 depth-2/3/4 质量前沿，提供质量模式与 cost-aware 模式。 & 独立 seed 泛化和 $n=21,22,23,24$ bridge。 \\
Guard & 用 baseline-preserving 规则兜底。 & 无 score-loss guard 与运行时节省。 \\
\bottomrule
\end{tabularx}
\end{table}

\section{实验设置}

实验回答三个问题。第一，\method{} 在完整 truth-table 可验证的小规模函数上是否确实降低 T-count 和 score。第二，在高维函数和项集合上，结构筛选、frontier policy 和 guard 是否仍能运行并保持语义正确。第三，与外部逻辑网络或 phase/Rz 工具链相比，本文优势是否仍成立，或暴露出哪些 gate-set 边界。

传统函数切片包含 $n\leq6$ 的 177 个函数，baseline 包括 direct \anf{}、AND-direct \anf{}、ESOP cube beam、weighted ESOP-MILP、SSHR-H、ABC/BDD 估计、ROS-style LUT proxy、mockturtle KLUT-to-XAG probe 和 RevKit Python \texttt{oracle\_synth}。高维切片包含 $n=14,15,16,18$ 的外部导出函数、$n=18,20$ 函数级边界、$n=20$--$40$ 项集级 scale，以及 $n=21,22,23,24$ truth-table bridge。Paired comparison 只纳入函数名或项集名匹配、语义正确且未 timeout 的结果。Timeout 行作为可扩展性边界保留，但不进入资源均值。

\begin{table}[H]
\centering
\scriptsize
\caption{实验层次与回答的问题。}
\label{tab:experiment-layers}
\begin{tabularx}{\linewidth}{>{\raggedright\arraybackslash}p{0.23\linewidth}>{\raggedright\arraybackslash}p{0.31\linewidth}>{\raggedright\arraybackslash}X}
\toprule
实验层次 & 数据与对照 & 回答的问题 \\
\midrule
传统函数 & $n\leq6$，177 个函数；ESOP、SSHR、ABC/BDD、mockturtle probe、RevKit API。 & 小规模完整验证条件下是否形成低 T-count 与低 score 优势。 \\
ROS-style LUT proxy & ABC $K=3,4,5$ sweep，309 个函数，927 行检查。 & 更强 LUT proxy 是否削弱本文优势。 \\
mockturtle XAG probe & ABC $K=4$ BLIF 到官方 mockturtle KLUT-to-XAG resynthesis。 & official-header probe 是否可运行，以及本文是否仍优于该 XAG proxy。 \\
高维函数边界 & $n=18,20$ 函数级切片。 & 高维完整 truth-table 条件下哪些搜索分支仍有效。 \\
项集级 scale & $n=20$--$40$，6600 个方法行。 & 结构策略是否可扩展，以及 plan 和 emitted circuit 是否符号等价。 \\
Truth-table bridge & $n=21,22,23$ 的 300 个方法行，加 $n=24$ 的 60 个方法行完整验证。 & 项集级符号验证和函数级 oracle 验证是否一致，以及完整验证边界能否继续上移。 \\
\bottomrule
\end{tabularx}
\end{table}

\section{结果}

\subsection{传统函数上的资源优势}

表~\ref{tab:traditional} 给出 $n\leq6$ 传统函数上的平均逻辑资源。Pareto-\method{} 相对 direct \anf{} 平均 T-count 降低 72.25\%，平均 score 降低 67.80\%；相对 ESOP cube beam 获得 174/0/3 的 score 胜/负/平，平均 score 降低 36.09\%；相对 weighted ESOP-MILP 获得 167/3/7，平均 score 降低 29.84\%。SSHR-H 的平均 CNOT 最低，因此本文不能声称 CNOT-only 支配 SSHR。更稳妥的结论是，本文用一定 CNOT 和辅助比特 trade-off 换取更低 T-count、depth 和加权 score。

\begin{table}[H]
\centering
\small
\caption{$n\leq6$ 传统函数切片上的平均逻辑资源。}
\label{tab:traditional}
\begin{tabular}{lrrrrr}
\toprule
方法 & 平均 T & 平均 CNOT & 平均 depth & 平均 ancilla & 平均 score \\
\midrule
Direct ANF & 184.1 & 134.2 & 134.6 & 1.33 & 194.3 \\
AND-direct ANF & 101.0 & 166.5 & 166.9 & 2.18 & 115.2 \\
\method{} & 43.9 & 89.9 & 94.5 & 1.92 & 53.2 \\
Pareto-\method{} & \textbf{40.4} & 83.0 & \textbf{88.0} & 2.03 & \textbf{49.6} \\
ESOP-MILP & 83.6 & 133.5 & 159.6 & 2.32 & 96.7 \\
SSHR-H & 81.0 & \textbf{67.6} & 88.7 & 1.36 & 88.2 \\
\bottomrule
\end{tabular}
\end{table}

\subsection{外部逻辑网络 probe}

ROS-style LUT proxy 先用 ABC 对导出 BLIF 运行 $K=3,4,5$ sweep，再把每个映射后的 LUT 网络按 local \anf{} compute/action/uncompute 方式估计逻辑资源。该 proxy 覆盖 309 个函数，共 927 个映射行，全部通过 \ttable{} 检查。Best-$K$ proxy 相对固定 $K=4$ 自身获得 219/0/90 的 score 胜/负/平，平均 score 降低 18.12\%，说明它确实比单一 $K$ 设置更强。在这个更强 baseline 上，\method{} 和 Pareto-\method{} 均在 309/309 个函数上取得更低 score，平均 score 分别降低 83.77\% 和 84.27\%。

mockturtle probe 进一步把外部比较从 ABC-LUT 内部 proxy 推进到官方 mockturtle 头文件生态。流程为：先用 ABC 将导出 BLIF 映射为 $K=4$ LUT 网络，再用 mockturtle \texttt{blif\_reader} 读入 KLUT network，并调用 \texttt{xag\_npn\_resynthesis} 重综合为 XAG。该流程不是官方 ROS，不包含 SAT garbage management，也不输出可逆 Clifford+T 线路；它的作用是给出一个可运行、可编译、可统计的 official-header KLUT-to-XAG probe。

\begin{table}[H]
\centering
\scriptsize
\caption{官方 mockturtle KLUT-to-XAG probe 对比。资源为 XAG AND/XOR/depth 派生的逻辑层 proxy，不是官方 ROS、可逆 garbage management 或硬件 mapping。}
\label{tab:mockturtle}
\begin{tabular}{lrrr}
\toprule
Target vs mockturtle XAG K4 & Items & W/L/T & Mean $\Delta$ \\
\midrule
and\_resource\_nmcts & 177 & 165/12/0 & -28.97\% \\
and\_pareto\_resource\_nmcts & 177 & 166/11/0 & -31.50\% \\
and\_fprm\_polarity\_archive & 177 & 156/21/0 & -25.11\% \\
and\_direct\_anf & 177 & 19/158/0 & +50.20\% \\
direct\_anf & 177 & 9/168/0 & +142.26\% \\
and\_esop\_milp & 177 & 92/85/0 & +8.18\% \\
sshr\_h & 177 & 50/127/0 & +33.30\% \\
\bottomrule
\end{tabular}


\vspace{0.8em}
\begin{tabular}{lrrr}
\toprule
Target vs mockturtle XAG K4 & Items & W/L/T & Mean $\Delta$ \\
\midrule
and\_resource\_nmcts & 64 & 64/0/0 & -91.41\% \\
and\_profile\_resource\_nmcts & 64 & 64/0/0 & -91.41\% \\
and\_pareto\_resource\_nmcts & 64 & 64/0/0 & -91.49\% \\
and\_direct\_anf & 64 & 64/0/0 & -88.40\% \\
direct\_anf & 64 & 64/0/0 & -82.58\% \\
\bottomrule
\end{tabular}

\end{table}

表~\ref{tab:mockturtle} 显示，在 $n\leq6$ 的 177 个传统函数上，mockturtle probe 177/177 行正确、无 error 或 timeout；Pareto-\method{} 相对该 probe 为 166/11/0，平均 score 降低 31.50\%。在 $n=14$ 的 64 个高维函数上，probe 64/64 行正确，Pareto-\method{} 为 64/0/0，平均 score 降低 91.49\%。需要注意的是，高维切片中本文方法的 depth proxy 可能高于 XAG probe；主要优势来自 T、CNOT 和 peak ancilla 的大幅降低。因此这里的主张仍是加权逻辑资源优势，而不是所有分项无条件支配。

\subsection{RevKit phase/Rz 边界与内部 emitter}

RevKit Python \texttt{oracle\_synth} 在 $n\leq6$ 的 177 个函数上全部返回，但返回的是 Rz-phase netlist，而不是可直接按 T/Tdg 完整计数的 Clifford+T netlist。角度审计显示，只有 6 行不含非 Clifford \rz{}；171 行含非 Clifford \rz{}，总数为 9242，角度除以 $\pi$ 的最大分母为 64。因此，\method{} 相对 RevKit lower-bound score 的失败不能解释为精确 Clifford+T T-count 失败，而应解释为不同 gate-set/phase 口径下的压力测试。

\begin{table}[H]
\centering
\small
\caption{RevKit phase/Rz 边界结果。}
\label{tab:rz-boundary}
\begin{tabularx}{\linewidth}{>{\raggedright\arraybackslash}Xrr}
\toprule
比较口径 & 胜/负/平 & 平均相对变化 \\
\midrule
\method{} vs RevKit lower-bound score & 6/171/0 & $+751.69\%$ \\
\method{} vs RevKit score+1/Rz & 140/37/0 & $-14.52\%$ \\
Resource-NMCTS family vs RevKit score+1/Rz & 157/20/0 & $-17.89\%$ \\
Resource-NMCTS family vs RevKit score+1.5/Rz & 177/0/0 & $-42.26\%$ \\
Resource-NMCTS family vs RevKit $T/R_z=30$ proxy & 177/0/0 & $-95.03\%$ \\
\bottomrule
\end{tabularx}
\end{table}

为了避免 phase/Rz 讨论只停留在 RevKit 侧的成本敏感性，本文进一步实现了一个内部 phase-parity \anf{} emitter。对每个非空 \anf{} 单项式 $S$，它使用恒等式
\begin{equation}
  \prod_{i\in S}x_i =
  \frac{1}{2^{|S|-1}}\sum_{\emptyset\neq T\subseteq S}(-1)^{|T|+1}
  \bigoplus_{i\in T}x_i
  \label{eq:phase-parity}
\end{equation}
把单项式相位展开为 parity-phase gadgets，并合并相同 parity mask 的旋转角。常数项被视为全局相位，因此验证目标是 phase oracle up to global phase。该 emitter 在 $n\leq6$ 的 177 个传统函数上全部通过 Fraction 精确验证，平均 lower-bound score 为 10.11，平均 CNOT 为 87.40，平均 depth 为 114.96，平均 \rz{} 总数为 27.56，平均非 Clifford \rz{} 为 21.75，最大角度分母为 32。

\begin{table}[H]
\centering
\small
\caption{Phase-parity \anf{} emitter 与 RevKit \texttt{oracle\_synth} 的 phase/Rz 同口径比较。该 emitter 实现 phase oracle up to global phase；lower-bound score 只计 T-like rotations，非 Clifford \rz{} 单独报告。}
\label{tab:phase-parity}
\begin{tabular}{lrrrr}
\toprule
Metric & Items & W/L/T & Mean $\Delta$ \\
\midrule
lower-bound score & 177 & 40/137/0 & +69.25\% \\
score+1/Rz & 177 & 177/0/0 & -48.16\% \\
score+1.5/Rz & 177 & 177/0/0 & -53.26\% \\
score+2/Rz & 177 & 177/0/0 & -56.05\% \\
$T/R_z=30$ proxy & 177 & 177/0/0 & -64.98\% \\
non-Clifford Rz & 177 & 171/0/6 & -63.33\% \\
total Rz & 177 & 173/3/1 & -45.10\% \\
\bottomrule
\end{tabular}

\end{table}

表~\ref{tab:phase-parity} 的结论并不是“naive phase expansion 已经解决问题”。相反，它显示 naive parity-phase 展开在 RevKit lower-bound score 下只有 40/137/0，平均 score 反而高 69.25\%；但它的非 Clifford \rz{} 数量相对 RevKit 为 171/0/6，平均降低 63.33\%，总 \rz{} 数量为 173/3/1，平均降低 45.10\%。一旦每个非 Clifford \rz{} 加入 1 个 score 单位，它相对 RevKit 转为 177/0/0，平均 score 降低 48.16\%；在 $T/R_z=30$ 近似旋转综合 proxy 下同样为 177/0/0，平均降低 64.98\%。因此，phase-parity baseline 的价值在于把 phase/Rz 边界从“RevKit 的成本重解释”推进为“项目内部已有可验证 emitter”，同时也说明下一步必须做学习式或优化式 phase/Rz-aware search，以减少 parity gadget 数和旋转谱，而不是依赖朴素展开。

为了让 phase/Rz 方向真正具有搜索结构，本文进一步加入 fixed-polarity phase-parity search。对每个极性 $p$，先构造 $g(z)=f(z\xor p)$，计算 $g$ 的 \anf{}，再把每个 shifted parity mask 翻译回 $x$ 变量。若 $p$ 与 parity mask 的交叠为奇数，则
\begin{equation}
  \bigoplus_{i\in T} z_i
  =
  1 -
  \bigoplus_{i\in T} x_i ,
\end{equation}
因此对应旋转角取相反数，并把原角度并入全局相位。验证时不再只检查常数项全局相位，而是检查所有输入上的 phase difference 是否为同一个有理数。这个检查更适合 phase oracle，因为极性平移会产生非整数的全局相位。

\begin{table}[H]
\centering
\small
\caption{Fixed-polarity phase-parity search。每个函数枚举全部 $2^n$ 个极性，按 lower-bound score、score+1/Rz 和 $T/R_z=30$ proxy 三个目标分别选择最优极性；共 531 行全部 up-to-global-phase 验证通过。}
\label{tab:phase-parity-fprm}
\begin{tabular}{lrrr}
\toprule
Comparison & Items & W/L/T & Mean $\Delta$ \\
\midrule
FPRM-score vs ANF & 177 & 59/0/118 & -3.98\% \\
FPRM-Rz1 vs ANF & 177 & 59/0/118 & -1.65\% \\
FPRM-$T/R_z=30$ vs ANF & 177 & 59/0/118 & -0.47\% \\
FPRM-$T/R_z=30$ vs RevKit & 177 & 177/0/0 & -65.16\% \\
FPRM non-Clifford Rz vs RevKit & 177 & 171/0/6 & -63.48\% \\
FPRM CNOT vs RevKit & 177 & 34/140/3 & +34.18\% \\
\bottomrule
\end{tabular}

\end{table}

表~\ref{tab:phase-parity-fprm} 显示，\fprm{} 极性搜索相对原始 phase-parity \anf{} 没有产生负退化：按 lower-bound score 选择时为 59/0/118、平均降低 3.98\%；按 score+1/Rz 选择时为 59/0/118、平均降低 1.65\%；按 $T/R_z=30$ proxy 选择时为 59/0/118、平均降低 0.47\%。在 $T/R_z=30$ 目标下，59/177 个函数选择非零极性，平均极性权重为 0.57，平均非 Clifford \rz{} 从 21.75 降到 21.60。这个增量是正向的，但幅度仍小；它说明极性搜索是合理的 phase/Rz 搜索维度，而不是已经形成足够强的 phase/Rz-aware 方法。

\subsection{高维结构策略与验证桥接}

在 $n=18$ 的 12 个随机 \anf{} 函数上，adaptive depth-2 Boolean-ring screen 相对 single screen 获得 11/0/1，平均 score 从 11349.06 降至 10670.94；完整 \method{} 平均 score 为 7315.62，相对 adaptive screen 进一步降低 23.85\%。在 $n=20$ 边界函数上，FPRM root beam 和 fast linear-pair 均触发 300 s task timeout；adaptive depth-2 screen 相对 single screen 获得 5/0/1，平均 score 降低 7.47\%；screen-gated \method{} 可把平均时间从 77.10 s 降至 20.71 s。这个结果说明高维中结构筛选是主要质量来源之一，Resource tail 的收益取决于切片和预算。

Depth-frontier policy 将高预算 depth-2/3/4 质量前沿学习化。在 $n=20,28,40$ scale harness 中，它相对 fixed depth-2 获得 35/0/37，平均 score 降低 2.19\%；相对 all-depth adaptive depth$\leq4$ 平均 score 只高 0.97\%，但节省 58.69\% 时间。扩大训练集后的 large frontier policy 在独立 seed $n=24,28,32,40$ 泛化集上相对 fixed depth-2 获得 56/0/40，平均 score 降低 2.34\%；相对 all-depth adaptive 只高 0.10\%，仍节省 53.50\% 时间。Cost-aware frontier policy 则以 +0.99\% score 代价换取相对 large policy 的 33.02\% plan time 下降和 7.61\% 辅助线 lifetime area 下降。

新增 $n=24$ bridge 使用式~\eqref{eq:truth-mask} 的位掩码 truth-table evaluator，在 6 个生成式 \anf{} 函数上产生 60 个方法行。所有方法行均通过完整 \ttable{} oracle 验证、plan \anf{} 展开验证和 emitted-circuit 符号验证，mismatch 为 0。表~\ref{tab:n24-bridge} 显示，在该完整枚举切片上，adaptive all-depth 相对 single screen 为 6/0/0、平均 score 降低 9.42\%；相对 fixed depth-2 为 4/0/2、平均 score 降低 2.21\%。Depth-frontier policy 相对 fixed depth-2 为 3/0/3、平均 score 降低 2.05\%；相对 depth-4 只高 0.16\%，但 plan time 降低 27.08\%。该结果把完整验证边界从 $n=23$ 推到 $n=24$，同时延续了 frontier policy 的质量-时间折中结论。

\begin{table}[H]
\centering
\scriptsize
\caption{$n=24$ 完整 truth-table bridge 的资源比较。}
\label{tab:n24-bridge}
\begin{tabular}{rllrrr}
\toprule
n & Method & Baseline & Score W/L/T & Mean $\Delta$ score & Mean $\Delta$ plan time \\
\midrule
24 & adaptive\_all\_depth & screen\_single & 6/0/0 & -9.42\% & +34710.06\% \\
24 & adaptive\_all\_depth & screen\_depth2 & 4/0/2 & -2.21\% & +1070.36\% \\
24 & depth\_policy & screen\_single & 6/0/0 & -6.96\% & +2155.21\% \\
24 & depth\_frontier\_policy & screen\_depth2 & 3/0/3 & -2.05\% & +555.27\% \\
24 & depth\_frontier\_policy & screen\_depth4 & 0/1/5 & +0.16\% & -27.08\% \\
24 & depth2\_guard\_direct & screen\_depth2 & 0/0/6 & +0.00\% & +0.00\% \\
24 & screen\_depth3 & screen\_depth2 & 4/0/2 & -1.45\% & +194.81\% \\
24 & screen\_depth4 & screen\_depth2 & 4/0/2 & -2.21\% & +642.76\% \\
all & adaptive\_all\_depth & screen\_single & 6/0/0 & -9.42\% & +34710.06\% \\
all & adaptive\_all\_depth & screen\_depth2 & 4/0/2 & -2.21\% & +1070.36\% \\
all & depth\_policy & screen\_single & 6/0/0 & -6.96\% & +2155.21\% \\
all & depth\_frontier\_policy & screen\_depth2 & 3/0/3 & -2.05\% & +555.27\% \\
all & depth\_frontier\_policy & screen\_depth4 & 0/1/5 & +0.16\% & -27.08\% \\
all & depth2\_guard\_direct & screen\_depth2 & 0/0/6 & +0.00\% & +0.00\% \\
all & screen\_depth3 & screen\_depth2 & 4/0/2 & -1.45\% & +194.81\% \\
all & screen\_depth4 & screen\_depth2 & 4/0/2 & -2.21\% & +642.76\% \\
\bottomrule
\end{tabular}

\end{table}

\begin{table}[H]
\centering
\small
\caption{高维验证桥接与结构策略结果。}
\label{tab:highdim-validation}
\begin{tabularx}{\linewidth}{>{\raggedright\arraybackslash}p{0.36\linewidth}>{\raggedright\arraybackslash}p{0.24\linewidth}>{\raggedright\arraybackslash}X}
\toprule
验证项或比较 & 结果 & 说明 \\
\midrule
项集级 plan \anf{} 验证 & 6600/6600 & 覆盖 $n=20$ 到 $40$ 的 scale/frontier/frontier-policy rerun 与泛化 run。 \\
项集级 emitted-circuit 符号验证 & 6600/6600 & GF(2) 多项式模拟通过，mismatch 总数为 0。 \\
$n=21,22,23$ 完整 \ttable{} oracle 验证 & 300/300 & 所有 emitted circuit 逐点验证通过。 \\
$n=24$ 完整 \ttable{} oracle 验证 & 60/60 & 平均 truth-table 构造时间 0.18 s/函数；plan 与 emitted-circuit 符号验证也均为 60/60。 \\
Large frontier vs fixed depth-2 & 56/0/40，score $-2.34\%$ & 独立 seed 泛化集。 \\
Large frontier vs all-depth adaptive & score $+0.10\%$，time $-53.50\%$ & 质量接近全深度前沿但时间更低。 \\
Cost-aware frontier vs large frontier & score $+0.99\%$，time $-33.02\%$，aux area $-7.61\%$ & 快速质量折中模式。 \\
\bottomrule
\end{tabularx}
\end{table}

\section{讨论：什么算明显提升}

本文当前已经有一条独立于 SSHR 的主线：\anf{}/\fprm{} 项集合给出可验证状态，Boolean-ring screen 提供结构候选，神经先验和 \mcts{} 负责搜索控制，frontier policy 和 guard 处理质量-时间折中与风险控制。这个主线比“改 SSHR”更适合作为论文方向，因为它可以自然对比 SSHR、ESOP、CNOT、T-count、辅助比特、线路长度、深度和 phase/Rz 口径。

但这还不是最终投稿形态。当前最强证据是逻辑层 bit-flip oracle synthesis 的低 T-count 与低 score；最弱证据是 phase/Rz-aware search 和标准外部工具链复现。RevKit、phase-parity \anf{} 与 fixed-polarity phase-parity search 共同说明，gate-set 口径会显著改变结论。现在 phase/Rz 方向已经不只是成本敏感性分析：它有一个可验证 emitter 和一个可验证极性搜索分支；但后者相对 \anf{} 的平均收益仍小，说明下一步需要更强的搜索动作，而不是只枚举输入极性。$n=24$ bridge 则说明正确性验证能力可以继续上移，但它只是小样本完整枚举，不等同于 $n=25$--$40$ 全部完成 truth-table simulation。

因此，本文的“明显提升”目标应设置为三个可检查终点。第一，把 fixed-polarity phase search 扩展为学习式或优化式 phase/Rz-aware search，在 177 个传统函数上不仅保持不退化，还要显著降低 parity gadget 数、非 Clifford \rz{} 数或综合后 proxy score，并保持 177/177 up-to-global-phase 验证。第二，继续复现官方 ROS 或 legacy RevKit/CirKit 完整流程，使外部对照从 adapter/probe 扩展为标准工具链。第三，补一个最小后端映射层评估，哪怕只覆盖代表性函数和固定硬件拓扑，也要把 CNOT-depth、routing overhead、辅助比特 lifetime 与逻辑层 score 的关系讲清楚。

\begin{table}[H]
\centering
\small
\caption{当前主张、证据与边界。}
\label{tab:evidence-map}
\begin{tabularx}{\linewidth}{>{\raggedright\arraybackslash}p{0.27\linewidth}>{\raggedright\arraybackslash}X>{\raggedright\arraybackslash}p{0.25\linewidth}}
\toprule
主张 & 当前证据 & 边界 \\
\midrule
本文路线独立于 SSHR & 状态、动作和验证都定义在 \anf{}/\fprm{} 项集合上；SSHR 只作为 baseline。 & 不能声称 CNOT-only 支配 SSHR。 \\
低 T-count 与低 score 是主优势 & $n\leq6$ 传统函数相对 ANF、ESOP 和 weighted ESOP-MILP 有稳定优势。 & 资源权重仍是逻辑层人为设定。 \\
外部逻辑网络 probe 下仍有优势 & ROS-style LUT proxy 为 309/0/0；mockturtle traditional 为 166/11/0，高维 $n=14$ 为 64/0/0。 & 仍不是官方 ROS 或可逆线路输出。 \\
Phase/Rz 边界已有实际 emitter 与极性搜索 & RevKit 171/177 行含非 Clifford \rz{}；phase-parity \anf{} emitter 177/177 验证通过；fixed-polarity phase search 531/531 验证通过，$T/R_z=30$ 目标相对 \anf{} 为 59/0/118、相对 RevKit 为 177/0/0。 & 朴素 phase-parity lower-bound 只有 40/137/0；fixed-polarity 搜索均值收益仍小，需要学习式 phase/Rz search 和实际旋转综合输出。 \\
高维结果具有语义可信度 & 6600/6600 项集级符号验证和 360/360 bridge 完整验证通过，其中 $n=24$ 为 60/60。 & $n=25$--$40$ 未做完整 \ttable{} 枚举。 \\
\bottomrule
\end{tabularx}
\end{table}

\section{结论}

本文提出 \method{}，将量子布尔函数 oracle 综合建模为 \anf{}/\fprm{} 项集合上的资源约束搜索问题，并结合神经动作先验、\mcts{}、布尔环线性因子筛选、结构深度策略、depth-frontier policy 和保守 guard 生成逻辑层可逆线路。实验表明，该方法在传统函数上相对 ANF、ESOP、SSHR-H 和 weighted ESOP-MILP 获得稳定低 T-count 与低 score 优势，在 ROS-style LUT proxy 和官方 mockturtle KLUT-to-XAG probe 下仍保持优势，并通过项集级符号验证和 $n=21,22,23,24$ 完整 truth-table bridge 补强了高维正确性证据。

当前稿件最适合定位为“资源约束量子布尔函数综合的搜索方法论文”。它不是 SSHR 的附属改进，也不是泛泛使用 AI 做线路生成；它的核心贡献是把 oracle synthesis 的动作空间、搜索策略、资源模型和语义验证放在同一个 \anf{}/\fprm{} 框架中。v32 新增的 fixed-polarity phase-parity search 说明 phase/Rz 方向已经进入可验证搜索阶段，但目前只形成小幅改进；v33 新增的 $n=24$ bridge 则把高维完整 oracle 验证边界继续上移。下一步若能把更强的 learned phase/Rz search、标准外部工具链复现和最小后端映射评估补齐，本文就可以从阶段性方法证据推进到更强的投稿版本。

\begin{thebibliography}{99}

\bibitem{meuli2019multiplicative}
G.~Meuli, M.~Soeken, E.~Campbell, M.~Roetteler, and G.~De Micheli.
\newblock The role of multiplicative complexity in compiling low T-count oracle circuits.
\newblock In \emph{Proceedings of ICCAD}, 2019.

\bibitem{meuli2022xag}
G.~Meuli, M.~Soeken, and G.~De Micheli.
\newblock Xor-And-Inverter Graphs for quantum compilation.
\newblock \emph{npj Quantum Information}, 8:7, 2022.

\bibitem{meuli2020ros}
G.~Meuli, M.~Soeken, M.~Roetteler, and G.~De Micheli.
\newblock ROS: Resource-constrained oracle synthesis for quantum computers.
\newblock \emph{Electronic Proceedings in Theoretical Computer Science}, 318:119--130, 2020.

\bibitem{wille2009bdd}
R.~Wille and R.~Drechsler.
\newblock BDD-based synthesis of reversible logic for large functions.
\newblock In \emph{Proceedings of DAC}, pp.~270--275, 2009.

\bibitem{yu2025backend}
M.~Yu, A.~Tempia Calvino, M.~Soeken, and G.~De Micheli.
\newblock Back-end-aware fault-tolerant quantum oracle synthesis.
\newblock In \emph{Proceedings of ASP-DAC}, pp.~930--937, 2025.

\bibitem{zheng2025sshr}
Q.~Zheng, Y.~Xu, J.~Zhang, Z.~Su, and S.~Zheng.
\newblock CNOT oriented synthesis for small-scale Boolean functions using spatial structures of parallelotopes.
\newblock In \emph{Proceedings of ICCAD}, 2025.

\bibitem{tsaras2024shortcircuit}
D.~Tsaras, A.~Grosnit, L.~Chen, Z.~Xie, H.~Bou-Ammar, and M.~Yuan.
\newblock ShortCircuit: AlphaZero-driven circuit design.
\newblock \emph{arXiv:2408.09858}, 2024.

\bibitem{rietsch2024unitary}
S.~Rietsch, A.~Y. Dubey, C.~Ufrecht, M.~Periyasamy, A.~Plinge, C.~Mutschler, and D.~D. Scherer.
\newblock Unitary synthesis of Clifford+T circuits with reinforcement learning.
\newblock In \emph{IEEE International Conference on Quantum Computing and Engineering}, pp.~824--835, 2024.

\bibitem{riu2025rlzx}
J.~Riu, J.~Nogu{\'e}, G.~Vilaplana, A.~Garcia-Saez, and M.~P. Estarellas.
\newblock Reinforcement learning based quantum circuit optimization via ZX-calculus.
\newblock \emph{Quantum}, 9:1758, 2025.

\bibitem{machiya2026monteq}
M.~Machiya, M.~Menickelly, P.~Hovland, and J.~Liu.
\newblock MonteQ: A Monte Carlo tree search based quantum circuit synthesis framework.
\newblock \emph{arXiv:2604.19029}, 2026.

\end{thebibliography}

\end{document}
